% titel: Blatt 5, A5: Elementare Funktionen
% skript: Kap. 2 (elementare Funktionen, Monotonie, Grenzwerte)

% Teil I: Werkzeuge, kein Lösungsweg.

Zu gegebenen Vorgaben (Monotonie, Funktionswerte, Grenzverhalten) jeweils
ein elementares $f$ mit passendem Definitionsbereich finden.

\begin{werkzeugkasten}[Bauplan: Qualität kontrollieren mit einer \emph{Glocke}]
Um eine Funktion mit gewünschter Lage des Maximums und gewünschtem
Grenzwert-Niveau zu bauen, trennt man zwei Einflüsse:
\begin{itemize}
  \item \textbf{Lage von Maximum/Minimum:} Im Exponenten sitzt eine um $x_0$
        zentrierte \emph{Glocke} $-\tfrac12(x-x_0)^2$. Erklärt wird das
        durch die Ableitung: $f'(x)$ verschwindet am Glockenmaximum $x=x_0$
        und hat dort einen Vorzeichenwechsel (links steigend, rechts
        fallend) -- eine echte, \emph{isolierte} Extremstelle.
  \item \textbf{Niveau des Grenzwerts:} Eine \emph{additive Konstante}
        verschiebt den Graph vertikal und setzt das asymptotische Niveau.
\end{itemize}
Beispiel zu (a): $f(x)=1+9\mathrm e^{-(x-2)^2}$. Die Konstante $1$ ist das
Niveau für $|x|\to\infty$ (die Glocke klingt nach aussen ab), die $9$
ist die Höhe über dem Niveau, und $x=2$ das Zentrum des Maximums.

\textit{Skript:} Kap. 2 (Grenzwerte, Exponentialfunktion).
\end{werkzeugkasten}

\begin{werkzeugkasten}[Bauplan: Pole und asymptotisches Verhalten]
\begin{itemize}
  \item \textbf{Ein Pol bei $x=a$} mit Grenzwert $+\infty$ (bzw.\ $-\infty$)
        entsteht, wenn der Nenner gegen $0$ geht und der Zähler endlich
        bleibt: ein Term mit $(x-a)$ im Nenner. Das Vorzeichen des
        Zählers und die Seitenwahl steuern, ob es $\pm\infty$ ist.
  \item \textbf{Grenzwert $-\infty$ für $x\to\infty$} erreicht man etwa über
        $-\ln(x-1)$ oder $-x$: Alle diese Funktionen fallen nach links-unten
        \emph{weg}.
  \item \textbf{Der Definitionsbereich ist vorgegeben.} Eine Funktion mit
        $\ln(x-1)$ darf nur auf $R_{>1}$ stehen; ein Nenner $(x^2-1)$
        erzwingt $R\setminus\{-1,1\}$.
\end{itemize}
Welche Kombination gewählt wird, ergibt sich aus den Vorgaben; oft
kombiniert man einen \emph{Pol-Baustein} additiv mit einer linearen Funktion,
um zusätzlich vorgegebene Funktionswerte zu treffen.

\textit{Skript:} Kap. 2 (Definitionsbereich, Grenzwerte).
\end{werkzeugkasten}

\begin{achtung}[Typische Fallen]
\begin{itemize}
  \item \textbf{Jede Bedingung einzeln prüfen:} Monotonie \emph{und} Wert
        $f(2)=0$ \emph{und} beide Grenzwerte. Eine vergessene Bedingung ist
        der häufigste Fehler.
  \item \textbf{Grenzwerte nicht verwechseln:} \emph{Grenzwert $1$ für
        $\pm\infty$} heisst, beide einseitigen Grenzwerte existieren und
        sind gleich $1$; das ist nicht dasselbe wie \emph{Wert $1$ an einer
        Stelle}.
  \item \textbf{Definitionsbereich einhalten:} $\ln(x-1)$ ist nur auf
        $R_{>1}$ sinnvoll; einen Pol bestimmt der Nenner, nicht der Zähler.
  \item \textbf{Maximum \emph{sitzen lassen}:} \emph{Monoton wachsend für $x<2$,
        fallend für $x>2$} bedeutet \emph{Maximum bei $x=2$}, nicht bei
        irgendeiner Stelle.
\end{itemize}
\end{achtung}
